\clearpage
\begin{trapbox}
\textbf{\textcolor{CTrap}{易错 1：把“一个周期”说成“最小正周期”}}
\begin{center}\begin{tikzpicture}[font=\small,>=stealth,line join=round]
\draw[gray!40,rounded corners=3pt] (0,0) rectangle (14.2,8.3);\node at (7.1,7.700000000000001) {例6：两轴相距2，只保证4是周期，未必最小};
\begin{scope}[shift={(1.58,3.814705882352941)},x=1.2cm,y=1.352941176470588cm]\draw[->,gray] (-0.4,0)--(4.6,0) node[right]{$x$};\draw[->,gray] (0,-1.6)--(0,1.8) node[above]{$y$};\draw[blue,thick,solid,domain=-0.1:4.2,samples=180] plot(\x,{cos(deg(3*pi*\x/2))});\fill[orange!85!black] (0,1) circle(2.2pt);\fill[orange!85!black] (1.3333333333333333,1) circle(2.2pt);\fill[orange!85!black] (2.6666666666666665,1) circle(2.2pt);\fill[orange!85!black] (4,1) circle(2.2pt);\end{scope}\node[anchor=west,align=left,text width=5.0cm] at (8.4,4.0) {$f=\cos\frac{3\pi x}2$\\[9pt]轴 $x=0,2$\\[9pt]公式给出周期4\\[9pt]最小正周期 $\frac43$\\[9pt]4含三个完整重复};
\node[text width=12.6cm,align=center] at (7.1,.65) {具体曲线在公式给出的距离内已经重复三次；最小周期需要额外判断。};
\end{tikzpicture}\end{center}



\textbf{错误：}“有两个对称轴且间距为 $2$，所以最小正周期一定是 $4$。”

\textbf{正确：}只能确定\textbf{$4$ 是一个周期}，最小正周期可能更小。见例 6：$\cos(3\pi x/2)$ 的最小正周期为 $4/3$，但 $4$ 同样是周期。

\medskip

\end{trapbox}
\clearpage
\begin{trapbox}
\textbf{\textcolor{CTrap}{易错 2：两处都是中心，就直接用 $2d$}}
\begin{center}\begin{tikzpicture}[font=\small,>=stealth,line join=round]
\draw[gray!40,rounded corners=3pt] (0,0) rectangle (14.2,8.3);\node at (7.1,7.700000000000001) {例7：两个中心一高一低，直线只会斜着搬移};
\begin{scope}[shift={(2.225,2.540322580645161)},x=0.75cm,y=0.49462365591397844cm]\draw[->,gray] (-1.5,0)--(6.5,0) node[right]{$x$};\draw[->,gray] (0,-1.8)--(0,7.5) node[above]{$y$};\draw[blue,thick,solid,domain=-1:6,samples=180] plot(\x,{\x});\fill[orange!85!black] (0,0) circle(2.2pt);\fill[orange!85!black] (2,2) circle(2.2pt);\fill[orange!85!black] (1,1) circle(2.2pt);\fill[orange!85!black] (5,5) circle(2.2pt);\draw[orange!85!black,dashed] (1,1)--(5,5);\end{scope}\node[anchor=west,align=left,text width=5.0cm] at (8.4,4.0) {$f=x$\\[9pt]中心 $(0,0),(2,2)$\\[9pt]$x\to x+4$ 时\\[9pt]$f(x+4)=f(x)+4$\\[9pt]不保持原高度};
\node[text width=12.6cm,align=center] at (7.1,.65) {两次中心对称会同时改变横纵坐标；同高度条件不可省略。};
\end{tikzpicture}\end{center}



必须检查两个中心的\textbf{高度相同}。若条件分别是
\[
\begin{aligned}
 f(a+t)+f(a-t)&=2c_1,\\
 f(b+t)+f(b-t)&=2c_2,
\end{aligned}
\]
则重复前面的换元可得一个更一般的关系：
\[
\boxed{f\bigl(x+2(b-a)\bigr)=f(x)+2(c_2-c_1).}
\]

只有 $c_1=c_2$ 才直接得到周期。\textbf{眼前一亮的理解：}两个中心若一高一低，图像每向右挪一段，可能也同时向上或向下挪，成为斜着重复，而不是原高度重复。例如 $f(x)=x$。

\medskip

\end{trapbox}
\clearpage
\begin{trapbox}
\textbf{\textcolor{CTrap}{易错 3：轴点组合刚移动 $2d$ 就断言是周期}}
\begin{center}\begin{tikzpicture}[font=\small,>=stealth,line join=round]
\draw[gray!40,rounded corners=3pt] (0,0) rectangle (14.2,8.3);\node at (7.1,7.700000000000001) {例3：轴在0，中心 $(1,3)$；走2高度尚未还原};
\begin{scope}[shift={(2.174626865671642,1.9785714285714284)},x=0.8955223880597015cm,y=0.6571428571428571cm]\draw[->,gray] (-1.2,0)--(5.5,0) node[right]{$x$};\draw[->,gray] (0,-0.5)--(0,6.5) node[above]{$y$};\draw[blue,thick,solid,domain=-0.8:5,samples=180] plot(\x,{3+2*cos(deg(pi*\x/2))});\fill[orange!85!black] (0,5) circle(2.2pt);\fill[orange!85!black] (2,1) circle(2.2pt);\fill[orange!85!black] (4,5) circle(2.2pt);\fill[orange!85!black] (1,3) circle(2.2pt);\draw[violet,dashed] (0,-.3)--(0,6);\draw[gray,dashed] (-1,3)--(5.2,3);\end{scope}\node[anchor=west,align=left,text width=5.0cm] at (8.4,4.0) {$f=3+2\cos\frac{\pi x}2$\\[9pt]轴 $x=0$，中心 $(1,3)$\\[9pt]$f(0)=5,\ f(2)=1$\\[9pt]$f(4)=5$\\[9pt]$d=1,\ T=4$};
\node[text width=12.6cm,align=center] at (7.1,.65) {相对高度3一高一低；移动2只翻转，移动4才保证原样重复。};
\end{tikzpicture}\end{center}



轴点组合真正得到的是
\[
 f\bigl(x+2(b-a)\bigr)=2c-f(x),
\]
不一定等于 $f(x)$。必须再进行一次同样的移动，才有
$f(x+4(b-a))=f(x)$。

\textbf{特殊提醒：}如果碰巧 $f(x)\equiv c$ 是常数函数，那么 $2d$ 也能作为周期；但\textbf{一般情况下不能保证 $2d$ 是周期}。

\medskip

\end{trapbox}
\clearpage
\begin{trapbox}
\textbf{\textcolor{CTrap}{易错 4：忽略 $a\neq b$}}
\begin{center}\begin{tikzpicture}[font=\small,>=stealth,line join=round]
\draw[gray!40,rounded corners=3pt] (0,0) rectangle (14.2,8.3);\node at (7.1,7.700000000000001) {轴与中心处于同一横坐标时，函数只能恒为c};
\begin{scope}[shift={(3.242857142857143,2.068181818181818)},x=0.8571428571428571cm,y=0.8363636363636363cm]\draw[->,gray] (-2.5,0)--(4.5,0) node[right]{$x$};\draw[->,gray] (0,-0.5)--(0,5) node[above]{$y$};\draw[blue,thick,solid,domain=-2:4,samples=180] plot(\x,{3});\fill[orange!85!black] (1,3) circle(2.2pt);\draw[violet,dashed] (1,-.3)--(1,4.8);\end{scope}\node[anchor=west,align=left,text width=5.0cm] at (8.4,4.0) {图取 $a=1,c=3$\\[9pt]既关于 $x=1$ 对称\\[9pt]又关于 $(1,3)$ 对称\\[9pt]只能 $f\equiv3$\\[9pt]任意正数都是周期};
\node[text width=12.6cm,align=center] at (7.1,.65) {此时距离d为0，不能套公式写正周期0；常数函数没有最小正周期。};
\end{tikzpicture}\end{center}



$T$ 必须\textbf{大于 $0$}，不能把 $a=b$ 代进去写出 $T=0$。

\textbf{特别情形：}若轴与中心刚好同在横坐标 $a$，既有 $f(2a-x)=f(x)$，又有 $f(2a-x)=2c-f(x)$，两式相减得 $f(x)=c$，整个函数只能是常数函数；这与通常的轴点公式前提不同。

\medskip

\end{trapbox}
\clearpage
\begin{trapbox}
\textbf{\textcolor{CTrap}{易错 5：对称中心不是“关于 $x=b$ 的轴”}}
\begin{center}\begin{tikzpicture}[font=\small,>=stealth,line join=round]
\draw[gray!40,rounded corners=3pt] (0,0) rectangle (14.2,8.5);\node at (7.1,7.9) {轴对称保持同高，中心对称保持两值之和};
\node at (2.4,6.9) {$f=\cos\frac{\pi x}2$};\begin{scope}[shift={(2.4,4.341935483870968)},x=0.85cm,y=1.1612903225806452cm]\draw[->,gray] (-2,0)--(2,0) node[right]{$x$};\draw[->,gray] (0,-1.5)--(0,1.6) node[above]{$y$};\draw[blue,thick,solid,domain=-1.6:1.6,samples=160] plot(\x,{cos(deg(pi*\x/2))});\fill[orange!85!black] (-0.5,0.7071067811865476) circle(2pt);\fill[orange!85!black] (0.5,0.7071067811865476) circle(2pt);\end{scope}\node[align=center,text width=4cm] at (2.4,1.7) {轴 $x=0$\\两点同高 $\frac{\sqrt2}2$};\node at (7.050000000000001,6.9) {$g=\sin\frac{\pi x}2$};\begin{scope}[shift={(7.050000000000001,4.341935483870968)},x=0.85cm,y=1.1612903225806452cm]\draw[->,gray] (-2,0)--(2,0) node[right]{$x$};\draw[->,gray] (0,-1.5)--(0,1.6) node[above]{$y$};\draw[teal,thick,solid,domain=-1.6:1.6,samples=160] plot(\x,{sin(deg(pi*\x/2))});\fill[orange!85!black] (-1,-1) circle(2pt);\fill[orange!85!black] (1,1) circle(2pt);\end{scope}\node[align=center,text width=4cm] at (7.050000000000001,1.7) {中心 $(0,0)$\\$g(-1)+g(1)=0$};\node at (11.7,6.9) {$h\equiv0$};\begin{scope}[shift={(11.7,4.341935483870968)},x=0.85cm,y=1.1612903225806452cm]\draw[->,gray] (-2,0)--(2,0) node[right]{$x$};\draw[->,gray] (0,-1.5)--(0,1.6) node[above]{$y$};\draw[violet,thick,solid,domain=-1.6:1.6,samples=160] plot(\x,{0});\end{scope}\node[align=center,text width=4cm] at (11.7,1.7) {常数函数可以\\同时具有两种对称};
\node[text width=12.6cm,align=center] at (7.1,.65) {三图同尺度；相等与相加为常数不同，不能把一个中心读成一条轴。};
\end{tikzpicture}\end{center}



$ f(b+t)+f(b-t)=2c $ 表示\textbf{关于点 $(b,c)$ 对称}，左右两个函数值\textbf{相加为 $2c$}，不一定彼此相等。
只有 $f(b+t)=f(b-t)$ 才表示关于直线 $x=b$ 对称。

\medskip

\end{trapbox}
\clearpage
\begin{trapbox}
\textbf{\textcolor{CTrap}{易错 6：只在几处取值相同就说函数周期}}
\begin{center}\begin{tikzpicture}[font=\small,>=stealth,line join=round]
\draw[gray!40,rounded corners=3pt] (0,0) rectangle (14.2,8.3);\node at (7.1,7.700000000000001) {只找到两点同高，不能认定处处周期};
\begin{scope}[shift={(1.6901639344262296,2.033333333333333)},x=0.9836065573770493cm,y=0.3833333333333333cm]\draw[->,gray] (-0.6,0)--(5.5,0) node[right]{$x$};\draw[->,gray] (0,-1)--(0,11) node[above]{$y$};\draw[blue,thick,solid,domain=-0.2:5.1,samples=180] plot(\x,{(\x-2)^2});\fill[orange!85!black] (0,4) circle(2.2pt);\fill[orange!85!black] (4,4) circle(2.2pt);\fill[orange!85!black] (1,1) circle(2.2pt);\fill[orange!85!black] (5,9) circle(2.2pt);\end{scope}\node[anchor=west,align=left,text width=5.0cm] at (8.4,4.0) {$f=(x-2)^2$\\[9pt]$f(0)=f(4)=4$\\[9pt]但 $f(1)=1$\\[9pt]$f(5)=9\ne1$\\[9pt]4不是周期};
\node[text width=12.6cm,align=center] at (7.1,.65) {证明周期要求任意输入；一对相等值只说明那一对点。};
\end{tikzpicture}\end{center}



例如发现 $f(0)=f(4)$，并不能推断 $4$ 为周期。必须证明\textbf{对任意实数 $x$}，$f(x+4)=f(x)$ 都成立。

\medskip

\end{trapbox}
\clearpage
\begin{trapbox}
\textbf{\textcolor{CTrap}{易错 7：忘记定义范围}}
\begin{center}\begin{tikzpicture}[font=\small,>=stealth,line join=round]
\draw[gray!40,rounded corners=3pt] (0,0) rectangle (14.2,8.3);\node at (7.1,7.700000000000001) {截取一个波形，定义域无法按非零周期不断移动};
\begin{scope}[shift={(1.6625,3.814705882352941)},x=0.9375000000000001cm,y=1.352941176470588cm]\draw[->,gray] (-0.6,0)--(5.8,0) node[right]{$x$};\draw[->,gray] (0,-1.6)--(0,1.8) node[above]{$y$};\draw[blue,thick,solid,domain=0:4,samples=180] plot(\x,{cos(deg(pi*\x/2))});\fill[orange!85!black] (0,1) circle(2.2pt);\fill[orange!85!black] (4,1) circle(2.2pt);\end{scope}\node[anchor=west,align=left,text width=5.0cm] at (8.4,4.0) {$D=[0,4]$\\[9pt]$f=\cos\frac{\pi x}2$\\[9pt]看到一个完整波形\\[9pt]但 $f(5)$ 无定义\\[9pt]整个函数不能照搬周期4};
\node[text width=12.6cm,align=center] at (7.1,.65) {须检查整个定义域是否随周期平移保持不变；截取的片段不自动继承周期。};
\end{tikzpicture}\end{center}



本结论直接采用 $f$ 在整个实数范围有定义。若题目函数只在一部分 $x$ 上有定义，应先检查平移 $T$ 后是否还在定义域内，不能机械套用。
\end{trapbox}
